
The principal LQR comparison uses a separately validated control weight, rather than relying only on the reference design $Q=I_3$, $R=1$. A sensitivity test fixes $Q=I_3$ and considers $R\in\{0.01,0.1,1,10,100\}$. For each node, $R$ is selected using ten validation initial conditions and the same nominal, $\eta=1$, and impulse normalized-MAE score used for SMC. All five candidate scores are retained in the reproducibility records. Both nodes select $R=0.01$, which is then frozen before evaluating 20 held-out test initial conditions. This error-oriented validation does not penalize measured energy and has access to non-ideal validation scenarios, unlike nominal-only PPO training. The selected value lies at the lower edge of the tested grid and defines the validated LQR setting used here. The historical $R=1$ results and fixed gain $k=8$ are unchanged.

\begin{table}[pos=!t,width=\linewidth]
\centering
\small
\setlength{\tabcolsep}{2pt}
\caption{LQR weight sensitivity on held-out tests. Each entry is a mean over 20 initial conditions; PPO additionally averages three training seeds. The selected $R=0.01$ is frozen using separate validation seeds. MAE covers the full horizon, including the impulse scenario.}
\label{tab:lqr_weight_sensitivity}
\begin{tabular}{@{}clrrrrrr@{}}
\toprule
 & & \multicolumn{2}{c}{PPO} & \multicolumn{2}{c}{LQR, $R=1$} & \multicolumn{2}{c}{LQR, $R=0.01$} \\
Node & Scenario & MAE & $E_u$ & MAE & $E_u$ & MAE & $E_u$ \\
\midrule
2 & Nominal & 0.0930 & 9.09 & 0.0907 & 3.78 & 0.0772 & 10.58 \\
2 & $\eta=2$ & 0.1703 & 11.85 & 0.2330 & 5.58 & 0.1486 & 14.60 \\
2 & Impulse & 0.4339 & 17.51 & 0.5128 & 9.95 & 0.3368 & 31.23 \\
2 & $\sigma=0.05$ & 0.2669 & 135.57 & 0.3742 & 11.67 & 0.2427 & 90.27 \\
\midrule
3 & Nominal & 0.0431 & 6.54 & 0.0510 & 2.65 & 0.0392 & 4.99 \\
3 & $\eta=2$ & 0.0656 & 6.44 & 0.0989 & 3.56 & 0.0647 & 5.29 \\
3 & Impulse & 0.2581 & 18.09 & 0.3330 & 7.49 & 0.2592 & 13.71 \\
3 & $\sigma=0.05$ & 0.1312 & 69.96 & 0.1826 & 7.11 & 0.1264 & 23.27 \\
\bottomrule
\end{tabular}
\end{table}


Table~\ref{tab:lqr_weight_sensitivity} shows that the original LQR weighting materially affects the comparison. At Node 2, $R=0.01$ lowers MAE below PPO in all four tested settings. Under mismatch, its MAE and energy are 0.1486 and 14.60, compared with PPO's 0.1703 and 11.85. After the impulse, its mean post-pulse error is 1.5536 versus PPO's 2.0800, with energy 31.23 versus 17.51. Under noise, it improves both MAE (0.2427 versus 0.2669) and energy (90.27 versus 135.57). The selected weighting therefore changes the error ordering obtained with the reference $R=1$ design.

At Node 3, the selected LQR has lower mean error and energy than PPO in the nominal, mismatch, and noise tests. PPO has a slightly lower mean post-impulse error (1.2969 versus 1.3169), but higher energy (18.09 versus 13.71). The recovery means are close, with a clearer difference in energy. The result highlights controller weighting as an important factor in the observed accuracy--effort balance.
